% !PRINTED 80
\weberpage{80}

%<<CONT
ist, die Wurzel einer Abel'schen Gleichung ist, deren Gruppe mit der zu $\mathfrak{A}$ reciproken Gruppe $\mathfrak{B}$ isomorph ist.

Die Aufgabe, die wir jetzt stellen und lösen wollen, ist folgende:

\gesperrt{I. Es soll der Modul $m$ und die Gruppe $\mathfrak{A}$ auf alle mögliche Arten so bestimmt werden, dass $\mathfrak{B}$ ein beliebig gegebenes System von Invarianten hat.}

Oder was dasselbe ist:

\gesperrt{Es sollen alle Kreistheilungskörper von gegebener Gruppe bestimmt werden.}

Wir machen dabei immer die Voraussetzung, dass $\mathfrak{A}$ ein primärer Theiler von $\mathfrak{N}$ sein soll, weil wir sonst einen und denselben Körper mehrmals erhalten würden. Die Aufgabe zerfällt in zwei Theile, nämlich:

\gesperrt{II. Welche Moduln $m$ sind geeignet, eine Gruppe $\mathfrak{B}$ von gegebenen Invarianten zu erzeugen?}

\gesperrt{III. Wie findet man diese Gruppe $\mathfrak{B}$ und die zugehörige Gruppe $\mathfrak{A}$?}

Wir müssen bei dieser Untersuchung die Bezeichnung gegen das Frühere etwas ändern, damit die Uebersichtlichkeit nicht verloren geht. Die exceptionelle Stellung der Primzahl 2, die bei den bisherigen Betrachtungen immer berücksichtigt werden musste, machte eine etwas umständliche Bezeichnung nothwendig. Diese Unterscheidung ist im Folgenden nicht mehr in dem Maasse nothwendig, und darum lässt sich die Bezeichnung jetzt vereinfachen.

Wir bezeichnen die Indices einer Zahl $a$ aus der Gruppe $\mathfrak{A}$ mit

\[
\alpha_1, \alpha_2, \dots \alpha_\mu \quad \text{(Indices von $a$)}
\]

und die entsprechenden Indexmoduln mit

\[
c_1, c_2, \dots c_\mu \quad \text{(Indexmoduln)},
\]

so dass $\mu$ gleich der Anzahl der in $m$ aufgehenden Primzahlen, oder wenn $m$ durch 8 theilbar ist, um 1 grösser ist.

Es seien ferner

\[
\omega_1, \omega_2, \dots \omega_\mu
\]

primitive Einheitswurzeln der Grade $c_1, c_2, \dots c_\mu$. Sind nun die Indices einer Zahl $b$ in der zu $\mathfrak{A}$ reciproken Gruppe $\mathfrak{B}$

\[
\beta_1, \beta_2, \dots \beta_\mu \quad \text{(Indices von $b$),}
\]
%CONT>>
